\chapter{Infective Endocarditis}
\index{Infective Endocarditis}

\section*{Definition and Overview}
Infective endocarditis is an infection of the inner lining of the heart, most commonly of the heart valves. It occurs when bacteria or other organisms enter the bloodstream, settle on damaged or abnormal valves or on artificial valves, and form infected vegetations. These clumps of organisms and debris can damage the valve, cause heart failure, and break off to block blood vessels elsewhere in the body, causing stroke, abscesses, or kidney damage. Endocarditis is uncommon but serious, with significant mortality, and requires prolonged antibiotic treatment and often surgery. Prevention centres on good dental hygiene, prompt treatment of infections, and careful management of the small group of people at highest risk.

\section*{Epidemiology}
Infective endocarditis affects around three to ten people per 100,000 each year, and its incidence has increased with ageing populations, more prosthetic valves, and more people living with heart devices such as pacemakers. It is more common in men and in older people. Risk is highest in people with prosthetic valves, previous endocarditis, certain congenital heart conditions, and valve disease, while intravenous drug use is a major risk factor in younger people. Rheumatic valve disease remains an important cause in lower-income countries. Around one in four people with endocarditis has no identifiable underlying heart disease.

\section*{Aetiology and Pathophysiology}
Endocarditis usually begins with bacteraemia, the presence of bacteria in the blood. The mouth is a common source: dental procedures, poor oral hygiene, and gum disease allow oral streptococci to enter the bloodstream. Skin organisms enter through wounds, intravenous lines, and injecting drug use. In most people, the body clears transient bacteraemia rapidly, but on abnormal or prosthetic valves the organisms adhere, multiply, and form vegetations. The infection damages the valve leaflets, causes regurgitation and abscesses, and fragments of the vegetation can embolise to the brain, spleen, kidneys, or limbs. Staphylococcus aureus, streptococci, and enterococci cause most cases, with more resistant organisms in healthcare and drug-use settings.

\section*{Clinical Features}
\begin{itemize}
    \item Fever, often with chills and night sweats
    \item Fatigue, malaise, loss of appetite, and weight loss
    \item A new or changing heart murmur
    \item Small red spots in the skin, nail-splinter haemorrhages, and tender nodules on fingers or toes
    \item Embolic events: stroke, pain from blockage of a limb or organ, or kidney damage with blood in the urine
    \item Back or joint pain, common in some forms of endocarditis
    \item Heart failure with breathlessness and fluid retention as valves fail
    \item Symptoms may develop acutely or over weeks
\end{itemize}

\section*{Differential Diagnosis}
Endocarditis can be difficult to diagnose because its early features are non-specific.
\begin{itemize}
    \item \textbf{Influenza and other viral infections:} fever and malaise without a cardiac source
    \item \textbf{Pneumonia and urinary tract infection:} common infections with fever
    \item \textbf{Malaria and other tropical diseases:} fever in returned travellers
    \item \textbf{Autoimmune disease:} such as systemic lupus, which can mimic the features
    \item \textbf{Atrial myxoma:} a benign heart tumour with embolic features
    \item \textbf{Rheumatic fever:} valve involvement with a different mechanism
    \item \textbf{Other causes of sepsis:} bacteraemia without valve infection
\end{itemize}

\section*{When to Seek Medical Care}
Anyone with a fever, especially a persistent or unexplained fever, should be assessed, particularly if they have a prosthetic valve, previous endocarditis, congenital heart disease, or valve disease, or if they inject drugs. A new murmur with fever, or any signs of embolism such as a sudden neurological symptom or pain in a limb, demands urgent assessment. Good dental care and prompt treatment of infections reduce risk.

\textbf{Contact your local emergency services for:}
\begin{itemize}
    \item Sudden weakness, numbness, or difficulty speaking, which may indicate a stroke
    \item Severe breathlessness, chest pain, or collapse
    \item Sudden pain, pallor, or coldness of a limb, which may indicate embolism
    \item High fever with confusion or rapid deterioration
\end{itemize}

\section*{Diagnosis and Investigations}
\begin{itemize}
    \item \textbf{Blood cultures:} multiple sets taken before antibiotics, essential to identify the organism
    \item \textbf{Echocardiography:} transthoracic ultrasound first, with transoesophageal echocardiography for clearer views of valves and vegetations
    \item \textbf{Full blood count and inflammatory markers:} anaemia and elevated markers support the diagnosis
    \item \textbf{Urinalysis:} blood and protein indicating kidney involvement
    \item \textbf{Imaging:} CT, MRI, or PET to detect emboli and abscesses
    \item \textbf{Duke criteria:} the standard clinical framework combining microbiological and echocardiographic findings
    \item \textbf{Cardiac surgery review:} where complications develop or surgery is considered
\end{itemize}

\section*{Management}
\begin{itemize}
    \item \textbf{Intravenous antibiotics:} prolonged courses, usually four to six weeks in hospital, tailored to the organism and its sensitivities
    \item \textbf{Antibiotic stewardship:} selecting narrow-spectrum therapy once the organism is known
    \item \textbf{Surgery:} valve repair or replacement for heart failure, uncontrolled infection, abscesses, large vegetations, or recurrent embolism
    \item \textbf{Management of complications:} treatment of stroke, emboli, heart failure, and kidney injury
    \item \textbf{Treatment of drug dependence:} where injecting drug use is the cause
    \item \textbf{Follow-up:} long-term surveillance of valve function and the risk of recurrence
    \item \textbf{Preventive care after recovery:} dental care and advice on future procedures
\end{itemize}

\section*{Complications}
\begin{itemize}
    \item Heart failure from valve destruction
    \item Embolic stroke and infarction of other organs
    \item Abscesses in the heart or around the valve
    \item Kidney failure from emboli or immune damage
    \item Sepsis and septic shock
    \item Metastatic infection in bones, joints, or other sites
    \item Recurrence, especially in people who continue injecting drugs
    \item Death in around one in five to one in four cases, higher with staphylococcal infection
\end{itemize}

\section*{Prognosis and Outcomes}
Infective endocarditis remains a serious illness, with death rates around 15 to 25 per cent despite modern treatment. Outcomes are better when it is diagnosed early, the organism is identified, and complications are treated promptly with antibiotics and surgery where indicated. Valve repair or replacement restores many people to good function, and those who recover are advised on lifelong dental care and antibiotic prophylaxis before certain procedures. The best outcomes follow prevention: good oral hygiene, prompt treatment of infection, and careful care of lines and injection sites.

\section*{Prevention and Self-Care}
\begin{itemize}
    \item Maintain excellent oral hygiene: brush and floss daily and see a dentist regularly
    \item Treat gum disease and dental infections promptly
    \item Seek medical care early for persistent or unexplained fever
    \item Follow your doctor's advice on antibiotic prophylaxis before dental or surgical procedures if you have a high-risk heart condition
    \item Never inject drugs, and seek help for drug dependence
    \item Care for intravenous lines and wounds as instructed in hospital
    \item Attend follow-up appointments after treatment to monitor valve function
    \item Tell new doctors and dentists about any heart condition or prosthetic valve
\end{itemize}

\section*{Sources and Further Reading}
The following reputable sources provide further information for healthcare students and patients seeking reliable, current advice about infective endocarditis.
\begin{itemize}
    \item Heart Foundation Australia. \textit{Infective endocarditis information.} https://www.heartfoundation.org.au/
    \item Healthdirect Australia. \textit{Endocarditis.} https://www.healthdirect.gov.au/
    \item Australian Dental Association. \textit{Antibiotic prophylaxis and heart conditions.} https://www.ada.org.au/
    \item NHS. \textit{Endocarditis.} https://www.nhs.uk/conditions/endocarditis/
    \item Mayo Clinic. \textit{Endocarditis.} https://www.mayoclinic.org/
    \item MSD Manual. \textit{Infective endocarditis.} https://www.msdmanuals.com/
\end{itemize}
